\QHeading{第1题　口算。}
\par\noindent


\begin{oralgrid}
\OralItem{1}{350\times 2}{700}
\OralItem{2}{6\times 15}{90}
\OralItem{3}{4\times 15}{60}
\OralItem{4}{80\times 12}{960}
\OralItem{5}{25\times 8}{200}
\OralItem{6}{250\times 4}{1000}
\OralItem{7}{45\times 2}{90}
\OralItem{8}{2\times 25}{50}
\OralItem{9}{24\times 5}{120}
\OralItem{10}{150\times 4}{600}
\OralItem{11}{4\times 125}{500}
\OralItem{12}{80\times 25}{2000}\end{oralgrid}
\par

\QHeading{口算思路补充（对应上表题号）}

\Answer{上面列出得数，下面按小题给出不同计算思路。计数单位法适用于小数，整数题可比较凑整、拆分或商不变等方法。}
\MethodTitle{第2小题：$6\times 15$}

\MethodBlock{口算思路一：按数位拆分计算}{把一个因数拆成整十、整百数与余下部分之和。}{\begin{aligned}[t]&\phantom{={}}6\times 15\\
&=6\times (10+5)\\
&=6\times 10+6\times 5\\
&=60+30\\
&=90\end{aligned}}
\MethodTitle{第3小题：$4\times 15$}

\MethodBlock{口算思路一：按数位拆分计算}{把一个因数拆成整十、整百数与余下部分之和。}{\begin{aligned}[t]&\phantom{={}}4\times 15\\
&=4\times (10+5)\\
&=4\times 10+4\times 5\\
&=40+20\\
&=60\end{aligned}}
\MethodTitle{第4小题：$80\times 12$}

\MethodBlock{口算思路一：凑整后补偿}{把接近整十、整百的因数凑整，再补回或减去多算的部分。}{\begin{aligned}[t]&\phantom{={}}80\times 12\\
&=80\times (10+2)\\
&=80\times 10+80\times 2\\
&=800+160\\
&=960\end{aligned}}
\MethodTitle{第5小题：$25\times 8$}

\MethodBlock{口算思路一：拆因数，凑整计算}{把一个因数拆成两个因数，先凑成整百或整千。}{\begin{aligned}[t]&\phantom{={}}25\times 8\\
&=25\times (4\times 2)\\
&=(25\times 4)\times 2\\
&=100\times 2\\
&=200\end{aligned}}

\MethodBlock{口算思路二：凑整后补偿}{把接近整十、整百的因数凑整，再补回或减去多算的部分。}{\begin{aligned}[t]&\phantom{={}}25\times 8\\
&=25\times (10-2)\\
&=25\times 10-25\times 2\\
&=250-50\\
&=200\end{aligned}}
\MethodTitle{第7小题：$45\times 2$}

\MethodBlock{口算思路一：按数位拆分计算}{把一个因数拆成整十、整百数与余下部分之和。}{\begin{aligned}[t]&\phantom{={}}45\times 2\\
&=2\times (40+5)\\
&=2\times 40+2\times 5\\
&=80+10\\
&=90\end{aligned}}
\MethodTitle{第8小题：$2\times 25$}

\MethodBlock{口算思路一：按数位拆分计算}{把一个因数拆成整十、整百数与余下部分之和。}{\begin{aligned}[t]&\phantom{={}}2\times 25\\
&=2\times (20+5)\\
&=2\times 20+2\times 5\\
&=40+10\\
&=50\end{aligned}}
\MethodTitle{第9小题：$24\times 5$}

\MethodBlock{口算思路一：按数位拆分计算}{把一个因数拆成整十、整百数与余下部分之和。}{\begin{aligned}[t]&\phantom{={}}24\times 5\\
&=5\times (20+4)\\
&=5\times 20+5\times 4\\
&=100+20\\
&=120\end{aligned}}
\MethodTitle{第11小题：$4\times 125$}

\MethodBlock{口算思路一：按数位拆分计算}{把一个因数拆成整十、整百数与余下部分之和。}{\begin{aligned}[t]&\phantom{={}}4\times 125\\
&=4\times (100+25)\\
&=4\times 100+4\times 25\\
&=400+100\\
&=500\end{aligned}}
\MethodTitle{第12小题：$80\times 25$}

\MethodBlock{口算思路一：拆因数，凑整计算}{把一个因数拆成两个因数，先凑成整百或整千。}{\begin{aligned}[t]&\phantom{={}}80\times 25\\
&=25\times (4\times 20)\\
&=(25\times 4)\times 20\\
&=100\times 20\\
&=2000\end{aligned}}

\MethodBlock{口算思路二：按数位拆分计算}{把一个因数拆成整十、整百数与余下部分之和。}{\begin{aligned}[t]&\phantom{={}}80\times 25\\
&=80\times (20+5)\\
&=80\times 20+80\times 5\\
&=1600+400\\
&=2000\end{aligned}}


\par\noindent

\QHeading{第2题　用简便方法计算下面各题。}

\CalcBlock[方法一：简便计算]{1}{\begin{aligned}[t]&\phantom{={}}\color{black}101\times 46-46\\&\color{answerred}=(101-1)\times46\\&\color{answerred}=100\times46\\&\color{answerred}=4600\end{aligned}}
\par
\MethodBlock{方法二：按运算顺序直接计算（对照）}{本法只作结果核对；本题要求简便计算，前面的简便解法仍保留。}{\begin{aligned}[t]&\phantom{={}}101\times 46-46\\
&=4646-46\\
&=4600\end{aligned}}


\CalcBlock[方法一：简便计算]{2}{\begin{aligned}[t]&\phantom{={}}\color{black}153\times 16+47\times 16\\&\color{answerred}=(153+47)\times16\\&\color{answerred}=200\times16\\&\color{answerred}=3200\end{aligned}}
\MethodBlock{方法二：按运算顺序直接计算（对照）}{本法只作结果核对；本题要求简便计算，前面的简便解法仍保留。}{\begin{aligned}[t]&\phantom{={}}153\times 16+47\times 16\\
&=2448+752\\
&=3200\end{aligned}}

\CalcBlock[方法一：简便计算]{3}{\begin{aligned}[t]&\phantom{={}}\color{black}35\times 102\\&\color{answerred}=35\times(100+2)\\&\color{answerred}=35\times100+35\times2\\&\color{answerred}=3500+70\\&\color{answerred}=3570\end{aligned}}
\MethodBlock{方法二：按数位拆分计算}{把一个因数拆成整十、整百数与余下部分之和。}{\begin{aligned}[t]&\phantom{={}}35\times 102\\
&=102\times (30+5)\\
&=102\times 30+102\times 5\\
&=3060+510\\
&=3570\end{aligned}}

\CalcBlock[方法一：简便计算]{4}{\begin{aligned}[t]&\phantom{={}}\color{black}170\times 14-70\times 14\\&\color{answerred}=(170-70)\times14\\&\color{answerred}=100\times14\\&\color{answerred}=1400\end{aligned}}
\MethodBlock{方法二：按运算顺序直接计算（对照）}{本法只作结果核对；本题要求简便计算，前面的简便解法仍保留。}{\begin{aligned}[t]&\phantom{={}}170\times 14-70\times 14\\
&=2380-980\\
&=1400\end{aligned}}

\CalcBlock[方法一：简便计算]{5}{\begin{aligned}[t]&\phantom{={}}\color{black}32\times 99+32\\&\color{answerred}=32\times(99+1)\\&\color{answerred}=32\times100\\&\color{answerred}=3200\end{aligned}}
\MethodBlock{方法二：按运算顺序直接计算（对照）}{本法只作结果核对；本题要求简便计算，前面的简便解法仍保留。}{\begin{aligned}[t]&\phantom{={}}32\times 99+32\\
&=3168+32\\
&=3200\end{aligned}}

\CalcBlock[方法一：简便计算]{6}{\begin{aligned}[t]&\phantom{={}}\color{black}125\times 88\\&\color{answerred}=125\times(80+8)\\&\color{answerred}=125\times80+125\times8\\&\color{answerred}=10000+1000\\&\color{answerred}=11000\end{aligned}}
\MethodBlock{方法二：拆因数，凑整计算}{把一个因数拆成两个因数，先凑成整百或整千。}{\begin{aligned}[t]&\phantom{={}}125\times 88\\
&=125\times (8\times 11)\\
&=(125\times 8)\times 11\\
&=1000\times 11\\
&=11000\end{aligned}}

\MethodBlock{方法三：把特殊因数改成整百或整千}{先把25、125等改写成整百或整千除以一个数。}{\begin{aligned}[t]&\phantom{={}}125\times 88\\
&=(1000\div 8)\times 88\\
&=1000\times (88\div 8)\\
&=1000\times 11\\
&=11000\end{aligned}}

\MethodBlock{方法四：凑整后补偿}{把接近整十、整百的因数凑整，再补回或减去多算的部分。}{\begin{aligned}[t]&\phantom{={}}125\times 88\\
&=125\times (90-2)\\
&=125\times 90-125\times 2\\
&=11250-250\\
&=11000\end{aligned}}

\par\noindent

\QHeading{第3题　解决问题。}

\Problem{① 甲、乙两车同时从某地向相反方向沿直线行驶，甲车每小时行驶55千米，乙车每小时行驶45千米。5小时后，两车相距多少千米?}
\par


\MethodTitle{方法一}
\Answer{\Formula{\begin{aligned}[t]&\phantom{={}}(55+45)\times5\\&=100\times5\\&=500\text{（千米）}\end{aligned}}}
\Answer{答：两车相距500千米。}

\MethodTitle{方法二：分别计算两部分，再求和}
\Answer{不先合并每份的数量，而是分别求出两部分的总数，再相加。}
\Answer{\Formula{\begin{aligned}[t]&\phantom{={}}55\times 5+45\times 5\\
&=275+225\\
&=500\text{（千米）}\end{aligned}}}
\Answer{答：两车相距500千米。}
\Problem{② 水果店运来苹果和梨各25箱，每箱苹果重12千克，每箱梨重8千克。运来的水果一共重多少千克?}

\MethodTitle{方法一}
\Answer{\Formula{\begin{aligned}[t]&\phantom{={}}(12+8)\times25\\&=20\times25\\&=500\text{（千克）}\end{aligned}}}
\Answer{答：水果共重500千克。}

\MethodTitle{方法二：分别计算两部分，再求和}
\Answer{不先合并每份的数量，而是分别求出两部分的总数，再相加。}
\Answer{\Formula{\begin{aligned}[t]&\phantom{={}}12\times 25+8\times 25\\
&=300+200\\
&=500\text{（千克）}\end{aligned}}}
\Answer{答：水果一共重500千克。}
